The fate of a finite disturbance in a spatial system can depend on both its initial geometry and the environment through which it propagates. Here we investigate this dependence in a stochastic spatial evolutionary game coupled to a renewable, diffusing resource. Compact and dispersed perturbations containing the same number of defectors and the same resource deficit exhibit markedly different finite-time outcomes: dispersed perturbations are more likely to establish growth. Exact first-visit relations link this contrast to interfacial activity and the directional bias of strategy updates. As disturbances expand, they develop resource-depleted interiors surrounded by interfacial resource layers. Translating the resource field preserves the strategy configuration, resource distribution, and resource autocorrelation, yet changes the probability of subsequent collapse. Short and long translations produce responses of opposite sign. Exact resource-balance and sensitivity relations connect these effects to strategy--resource covariance and local payoff differences. Redistributing resources along a fixed interface can also reverse the initial growth response, although this ordering need not persist during subsequent evolution. These findings demonstrate that relative strategy--resource alignment contains information about finite-disturbance robustness that is not captured by the statistics of either field separately.
